\documentclass[11pt]{article}
\usepackage[margin=1in]{geometry}
\usepackage{amsmath,amssymb,amsthm}
\usepackage[T1]{fontenc}
\usepackage{lmodern}
\usepackage[hidelinks]{hyperref}
\newtheorem{theorem}{Theorem}
\newtheorem{lemma}{Lemma}
\newcommand{\PP}{\mathbb P}
\newcommand{\CC}{\mathbb C}
\newcommand{\eps}{\varepsilon}
\title{A Minimal-Surface Counterexample to a Seshadri Lower Bound}
\author{Alec Kriebel\\\small Independent researcher\\
\small\href{https://orcid.org/0009-0001-9320-500X}{ORCID: 0009-0001-9320-500X}}
\date{October 14, 2026}
\begin{document}
\maketitle
\begin{abstract}
We give an explicit family of counterexamples to the proposed lower bound
$\eps(S,L;x)\ge (2+|K_S^2|^{1/4})^{-1}$ for minimal complex projective
surfaces. For every integer $d\ge7$, a classical plane-pencil construction
followed by a double base change gives a smooth surface with ample canonical
bundle, $K_S^2=6(d-1)(d-3)$, and an ample integral line bundle $L$ with
$L^2=6$ and $\eps(S,L;x)\le1/(d-1)$. Already $d=7$ contradicts the bound.
The pencil and covering ingredients are classical; historical priority has
not been fully cleared. This unrefereed note used AI tools extensively and
has not received human peer review.
\end{abstract}

\section{The formulation and the result}
For a smooth complex projective surface $S$, an ample integral line bundle
$L$, and a closed point $x$, put
\[
 \eps(S,L;x)=\inf_{\substack{C\subset S\text{ integral curve}\\x\in C}}
                 \frac{L\cdot C}{\operatorname{mult}_x C}.
\]
Szemberg's final question in the 2007 preprint of \cite{Szemberg} asks for a
minimal polarized surface violating
\begin{equation}\label{target}
 \eps(S,L;x)\ \ge\ \frac{1}{2+|K_S^2|^{1/4}}.
\end{equation}
It is restated as Conjecture~2 in his Oberwolfach contribution
\cite[p.~1124]{OWR} and as Question~6.1.6 in \cite[p.~17]{Primer}, using the
cited preprint version for that locator. The point is arbitrary, and there
is no Picard-rank restriction. The additive $2$ lies outside the fourth
root; the problem catalogue's different transcription is not the primary
formula.

\begin{theorem}\label{main}
For every integer $d\ge7$ there are a smooth integral complex projective
surface $S$, an ample integral line bundle $L$ on $S$, and a point $x\in S$
such that $K_S$ is ample and
\[
 K_S^2=6(d-1)(d-3),\qquad L^2=6,\qquad
 \eps(S,L;x)\le\frac1{d-1}
       <\frac1{2+(K_S^2)^{1/4}}.
\]
In particular $S$ is minimal of general type and \eqref{target} is false.
\end{theorem}

The starting pencil and its polarization are Miranda's construction,
recorded by Ein--Lazarsfeld \cite[Example~3.1, p.~182]{EL}, Bauer
\cite[Proposition~3.3 and its proof]{Bauer}, and Bauer--Szemberg
\cite[Proposition~2.1, p.~2]{Adjoint}. We use a double base change branched
at four smooth fibres. We claim neither newness of these operations nor
historical firstness of this counterexample. Section~\ref{limits} gives the
priority and verification limits.

\section{An integral pencil and an ample line bundle}
We supply the existence argument at finite degree, rather than relying on
the asymptotic $d=m+1$ assertion in \cite[Proposition~3.3(b)]{Bauer}.
In $\PP^2$ use coordinates $[X:Y:Z]$ and set
\[
 f_d=Z(X^{d-1}-Y^{d-1})+X^d,\qquad C_0=V(f_d),\qquad p=[0:0:1].
\]
The coefficients of this linear polynomial in $Z$ are coprime in
$\CC[X,Y]$. Gauss's lemma, applied over $\CC(X,Y)$, proves irreducibility.
In the chart $Z=1$, the tangent cone at $p$ is $x^{d-1}-y^{d-1}$, a product
of distinct linear factors. Thus $p$ is an ordinary $(d-1)$-fold point.
It is the only singular point: the equations $f_Z=f_Y=0$ force either
$X=Y=0$, or $X,Y\ne0$ and $Z=0$; in the latter case $f_X=dX^{d-1}\ne0$.

Write $N_e=e(e+3)/2$. The reducible or nonreduced degree-$d$ locus $R$ in
$\PP^{N_d}$ is the finite union of the closed images of multiplication maps
\[
 \PP^{N_a}\times\PP^{N_{d-a}}\longrightarrow\PP^{N_d},
       \qquad 1\le a\le d-1.
\]
Since $N_d-N_a-N_{d-a}=a(d-a)\ge d-1$, its dimension is at most
$N_d-d+1$. The closure of the union of lines joining $[f_d]$ to $R$ has
dimension at most $N_d-d+2<N_d$ for $d\ge3$. Choose a form $g$ outside that
closure. Then the entire pencil $\langle f_d,g\rangle$ avoids $R$.
We can simultaneously choose $g$ smooth, nonzero at $p$, and transverse to
the smooth locus of $C_0$: these are nonempty open conditions in the same
irreducible parameter space, by Bertini and transversality for the complete
degree-$d$ system. B\'ezout gives exactly $d^2$ distinct simple base points
$p_i$, all away from $p$. Thus every pencil member is integral and reduced.

Blow up these base points once each, obtaining $b:Y\to\PP^2$. Let
$H=b^*\mathcal O_{\PP^2}(1)$, and let $E_i$ be the exceptional curves.
Transversality resolves the pencil and gives $f:Y\to\PP^1$, with fibre class
\[
 F=dH-\sum_{i=1}^{d^2}E_i.
\]
All fibres are strict transforms of integral reduced pencil members; the
$E_i$ are sections, not vertical components. General fibres are smooth.
The special fibre $C$, the strict transform of $C_0$, still has multiplicity
$d-1$ at the point $y$ above $p$.

Put $A=2F+E_1$. The blowup intersection rules give
\[
 F^2=0,\quad F\cdot E_i=1,\quad E_i^2=-1,\qquad
 A^2=3,\quad A\cdot E_1=1,\quad A\cdot F=1.
\]
Here $F=f^*\mathcal O_{\PP^1}(1)$ is nef. For every integral curve
$D\ne E_1$, $E_1\cdot D\ge0$. If $F\cdot D>0$, then $A\cdot D>0$.
If $F\cdot D=0$, $D$ is vertical, hence is an entire fibre because each
fibre is integral and reduced; in this case $A\cdot D=1$. Together with
$A\cdot E_1=1$ and $A^2>0$, the Nakai--Moishezon criterion proves that
$A$ is ample. The assertion about every fibre is essential here.

\section{The cover and minimality}
Choose four distinct smooth fibre values of $f$, avoiding the value of $C$.
Let $h:B\to\PP^1$ be the connected double cover simply branched at these
values. For example its equation in a coordinate avoiding the four values
at infinity is $v^2=\prod_{j=1}^4(t-t_j)$. Its smooth projective completion
is connected of genus one by Hurwitz. Form the fibre product
\[
 S=Y\times_{\PP^1}B,\qquad \pi:S\to Y,
\]
so $\pi$ is finite flat of degree two. It is \'etale off the four branch
fibres, including above the singularities of all other fibres. Along a
branch fibre, $f$ is smooth and the local cover has equation $t=v^2$ in
smooth coordinates $(u,t)$ on $Y$, so $S$ is smooth there as well.
The map $S\to B$ is proper with connected fibres, because they are the
fibres of $f$. Since $B$ is connected, $S$ is connected: a decomposition
into two open and closed subsets would give a decomposition of $B$, using
properness and connected fibres. A connected smooth scheme has a single
irreducible component here, since distinct components of a smooth scheme
cannot meet. Hence $S$ is integral. It is projective as a fibre product of
projective varieties.

The branch divisor on $Y$ is linearly equivalent to $4F$. The double cover
is associated with $\mathcal O_Y(2F)$. Hurwitz therefore gives the equality
of line bundles
\begin{equation}\label{canonical}
 K_S=\pi^*(K_Y+2F).
\end{equation}
This is the classical cyclic-cover formula; compare
\cite[Section~2, p.~3]{Fuentes}. In the present local model the ramification
divisor has line bundle $\pi^*\mathcal O_Y(2F)$ and is reduced, also giving
\eqref{canonical} directly. Since $K_Y=-3H+\sum E_i$, put
\[
 N=K_Y+2F=(d-3)H+F.
\]
We have $N\cdot E_i=1$, and for every other integral curve $D$, its image
in $\PP^2$ is a curve of degree $e>0$, so
\[
 N\cdot D=(d-3)e+F\cdot D>0\qquad(d\ge4).
\]
Moreover $H^2=1$ and $H\cdot F=d$, whence
\[
 N^2=(d-3)^2+2d(d-3)=3(d-1)(d-3)>0.
\]
Nakai--Moishezon proves $N$ ample. A finite surjective pullback of an ample
line bundle is ample (alternatively apply the projection formula and Nakai
on $S$). Thus \eqref{canonical} proves $K_S$ ample. In particular $S$ is of
general type and has no rational $(-1)$-curve, since adjunction would give
canonical degree $-1$ on such a curve. Consequently $S$ is minimal.

\section{The pointed curve and the comparison}
Set $L=\pi^*A$, an ample integral line bundle. The degree-two intersection
formula gives
\[
 L^2=2A^2=6,\qquad K_S^2=2N^2=6(d-1)(d-3).
\]
The fibre value of $C$ is unbranched. Its two inverse images under $h$ are
distinct complex points, so
\[
 \pi^{-1}(C)=C'\sqcup C'',\qquad C'\simeq C\simeq C''.
\]
In particular this is not a connected degree-two curve over $C$. Choose
$x\in C'$ mapping to $y$. The local isomorphism, or \'etaleness and the
split fibre identification, gives $\operatorname{mult}_x C'=d-1$.
Projection gives $L\cdot C'=A\cdot C=1$. Therefore
\[
 \eps(S,L;x)\le\frac{L\cdot C'}{\operatorname{mult}_x C'}=\frac1{d-1}.
\]
For $d\ge7$ the desired strict inequality is equivalent, after subtracting
$2$ and taking fourth powers of positive numbers, to
\[
 (d-3)^3>6(d-1).
\]
At $d=7$ this is $64>36$. The forward difference of the left side minus
the right side is $3(d-3)^2+3(d-3)-5>0$ for $d\ge7$, proving the claim
for every such integer. In the smallest example $K_S^2=144$, $L^2=6$,
and $144<4^4$ yields $1/6<(2+144^{1/4})^{-1}$. This proves
Theorem~\ref{main}.

Only the displayed upper bound is asserted; an exact Seshadri value is
unnecessary. The surfaces vary with $d$, and $x$ is a special point.
The construction does not resolve positivity of the infimum over all ample
integral line bundles and points on one fixed complex surface. It does not
contradict very-general-point or Picard-rank-one lower bounds, or canonical
and adjoint polarization bounds. For $d=4,5,6$ this particular curve does
not give the strict comparison; this says nothing about other curves or
lower bounds on those surfaces.

\section{Attribution, priority and verification limits}\label{limits}
Miranda's pencil polarization is classical, with the direct record in
\cite{EL}, its extension in \cite{Bauer}, and the close arbitrary-rational
construction in \cite{Adjoint}. Those examples use nonminimal blowups.
Earlier double-cover examples in \cite[Example~1.3]{Canonical} concern the
canonical polarization and produce a quotient $1/2$; that quotient cannot
violate \eqref{target}. Koll\'ar's example in \cite[Example~3.7]{EL}
concerns unbounded indices of very ampleness, which is a different assertion.
Fuentes Garc\'ia \cite{Fuentes} treats cyclic-cover constants in other
settings, including very general points over Picard-rank-one bases and
points on ramification divisors. The present curve is over an unbranched
singular fibre of a plane blowup.

Two independent internal primary-source priority audits, supplemented by
the author's source comparisons, found no explicit earlier resolution of
this precise bound in the text inspected. Historical priority nevertheless
cannot be fully cleared. In particular the full Bansal--Majumder 2025
survey \cite{Survey2025}, the full Roy 2019 thesis, and some other theses
were not inspected; longer sources and later citation bodies were partly
or selectively read, and final journal versions were not globally compared.
The verification package records exact reading scopes, access limits,
fingerprints and search ledgers. We make no claim of first discovery,
novelty of the covering method, exhaustive priority clearance, or independently
verified openness immediately before this note.

The accompanying authored-only verification package preserves the original
18 inputs from PR~585, new independent derivations and adversarial reviews,
source receipts, exact symbolic controls, and pinned nested inventories.
Its replay requires preinstalled SymPy~1.14.0 and mpmath~1.3.0; raw third-party
documents and account material are excluded. The finite controls supplement
the geometric proof, including its Bertini and Nakai--Moishezon inputs;
they do not replace existence arguments or establish priority. Integrity
checks do not certify authenticity or semantic authorship.

\paragraph{AI and review disclosure.}
AI tools were used extensively in solving, research, writing, source
comparison, reproduction and adversarial verification. Reviews described
here are internal AI reviews. This is an unrefereed preprint: it has not
received human peer review and is not certified by a proof assistant.

\begin{thebibliography}{10}
\raggedright
\bibitem{Szemberg} T.~Szemberg, \emph{An effective and sharp lower bound on
Seshadri constants on surfaces with Picard number~1}, J.~Algebra
319 (2008), 3112--3119.
\href{https://doi.org/10.1016/j.jalgebra.2007.10.036}{doi:10.1016/j.jalgebra.2007.10.036}.
Inspected version: \href{https://arxiv.org/abs/0711.0584v1}{arXiv:0711.0584v1}
(5 November 2007), final question, p.~7.
\bibitem{OWR} T.~Szemberg, \emph{Seshadri constants and geometry of surfaces},
in \emph{Multiplier Ideal Sheaves in Algebraic and Complex Geometry},
Oberwolfach Reports 21/2009, contribution pp.~1124--1126, Conjecture~2, p.~1124.
\href{https://doi.org/10.4171/OWR/2009/21}{doi:10.4171/OWR/2009/21}.
\bibitem{Primer} T.~Bauer, S.~Di Rocco, B.~Harbourne, M.~Kapustka,
A.~Knutsen, W.~Syzdek and T.~Szemberg, \emph{A primer on Seshadri constants},
Contemp.~Math. 496 (2009), 33--70.
\href{https://doi.org/10.1090/conm/496/09718}{doi:10.1090/conm/496/09718}.
Inspected version: \href{https://arxiv.org/abs/0810.0728v2}{arXiv:0810.0728v2}
(1 July 2010), Question~6.1.6, p.~17.
\bibitem{EL} L.~Ein and R.~Lazarsfeld, \emph{Seshadri constants on smooth
surfaces}, Ast\'erisque 218 (1993), 177--186.
\href{https://numdam.org/item/AST_1993__218__177_0/}{Numdam primary record}.
\bibitem{Bauer} T.~Bauer, \emph{Seshadri constants on algebraic surfaces},
Math.~Ann. 313 (1999), 547--583.
\href{https://doi.org/10.1007/s002080050272}{doi:10.1007/s002080050272}.
Inspected version: \href{https://arxiv.org/abs/math/9903072v1}{arXiv:math/9903072v1}.
\bibitem{Adjoint} T.~Bauer and T.~Szemberg, \emph{On the Seshadri constants
of adjoint line bundles}, \href{https://arxiv.org/abs/1011.0979v1}{arXiv:1011.0979v1}
(3 November 2010), Proposition~2.1, p.~2.
\bibitem{Fuentes} L.~Fuentes Garc\'ia, \emph{Cyclic coverings and Seshadri
constants on smooth surfaces},
\href{https://arxiv.org/abs/math/0511547v2}{arXiv:math/0511547v2}
(27 April 2006), Section~2, p.~3.
\bibitem{Canonical} T.~Bauer and T.~Szemberg, \emph{Seshadri constants on
surfaces of general type}, Manuscripta Math. 126 (2008), 167--175.
\href{https://doi.org/10.1007/s00229-008-0170-2}{doi:10.1007/s00229-008-0170-2}.
Inspected version: \href{https://arxiv.org/abs/0801.3245v1}{arXiv:0801.3245v1}.
\bibitem{Survey2025} A.~Bansal and S.~Majumder, \emph{Seshadri constants
and related conjectures over characteristic zero fields}, Bull.~Sci.~Math.
202 (2025), 103617.
\href{https://doi.org/10.1016/j.bulsci.2025.103617}{doi:10.1016/j.bulsci.2025.103617}.
Metadata inspected; full body unread and not used as proof or priority clearance.
\end{thebibliography}
\end{document}
